\documentclass[12pt]{article}
\usepackage[T1]{fontenc}
\usepackage[utf8]{inputenc}
\usepackage[brazil]{babel}
\usepackage[margin=1in]{geometry}
\usepackage{ragged2e}
\usepackage[normalem]{ulem}
\setlength{\parindent}{0pt}
\setlength{\parskip}{0.8\baselineskip}
\begin{document}
\justifying
\noindent Verifica-se que a questão objeto da controvérsia jurídica suscitada no recurso manejado pela parte recorrente esteve afetada no representativo da controvérsia - \textbf{Tema n. }\textbf{167} da TNU - (PU no processo n. PEDILEF 5003449-95.2016.4.04.7201/SC),  afetado em 13/09/2017 foi julgado (trânsito em julgado em 11/04/2018), por meio do qual foi elucidada a seguinte questão:

\noindent \textit{"Saber se o cálculo do salário de benefício do segurado que contribuiu em razão de atividades concomitantes vinculadas ao RGPS deve se dar com base na soma integral do salários de contribuição (respeitado o limite máximo) e sem a observância das limitações impostas pelo art. 32 da Lei 8.213/91."}

\noindent Restou firmada a seguinte tese:

\noindent \textit{"O cálculo do salário de benefício do segurado que contribuiu em razão de atividades concomitantes vinculadas ao RGPS e implementou os requisitos para concessão do benefício em data posterior a 01/04/2003, deve se dar com base na soma integral dos salários-de-contribuição (anteriores e posteriores a 04/2003) limitados ao teto. (Tese mantida, em face do julgamento do STJ no Tema 1070 no mesmo sentido)." (Tema 167 TNU).}

\noindent Com efeito, depreende-se da leitura do voto condutor do acórdão que o mesmo se encontra em conformidade com o decidido no  tema pela TNU.

\noindent A proposito, consta do acordao hostilizado: (...) RMI pelo art. 1º da Lei 6.423/77 - índices de atualização dos 24 1ºs salários-de-contribuição, anteriores aos 12 últimos, RMI - Renda Mensal Inicial, RMI - Renda Mensal Inicial, Reajustes e Revisões Específicas, DIREITO PREVIDENCIÁRIO, Competência, Jurisdição e Competência, DIREITO PROCESSUAL CIVIL E DO TRABALHO

\noindent Nessas condições, o conteúdo do Acórdão recorrido é consentâneo com o entendimento da TNU, firmado em julgamento de recurso representativo de controvérsia, o que atrai a incidência da regra prevista pelo art. 14, III, \textit{b,} da Resolução CJF n. 586/2019:

\noindent \textit{\textbf{Art. 14.}}\textit{ Decorrido o prazo para contrarrazões, os autos serão conclusos ao magistrado responsável pelo exame preliminar de admissibilidade, que deverá, de forma sucessiva: (...) }\textit{\textbf{III }}\textit{- negar seguimento a pedido de uniformização de interpretação de lei federal interposto contra acórdão que esteja em conformidade com entendimento consolidado: }\textit{\textbf{b)}}\textit{ em }\uline{\textit{\textbf{recurso representativo de controvérsia pela Turma Nacional de Uniformização}}}\textit{ ou em pedido de uniformização de interpretação de lei dirigido ao Superior Tribunal de Justiça;(...).}

\noindent Posto isso, com arrimo no \uline{\textbf{art. 14, III, }}\uline{\textit{\textbf{b}}}\textit{,} da Resolução CJF n. 586/2019 c/c art. 1.030, I, do CPC, \textbf{NEGO SEGUIMENTO }\textbf{a}o\textbf{ PU}\textbf{.}

\noindent Intimem-se. Aguarde-se o decurso do prazo recursal. Após, certifique-se o trânsito em julgado e devolva-se ao juízo de origem.
\end{document}
